% GENERATED FILE. Edit canonical lessons, then run npm run generate:books.
% canonical-source-hash: fnv1a64:518288635b22da65
% canonical-lessons: JA-C465-enogu, JA-C465-fude, JA-C465-egaku, JA-C465-gaka, JA-C465-bijutsu

\chapter{Paints and Brushes}
\label{ch:ja-paints-and-brushes}
\hlchaptermodality{465}

\begin{chapteropening}
\textbf{By the end of this chapter:} \emph{I can talk about painting.}

Complete the last lesson of chapter 465: I can talk about painting.
\end{chapteropening}

\section[\texorpdfstring{\ja{えのぐ} (enogu)}{えのぐ (enogu)}]{\ja{えのぐ} (enogu) — paints}

\label{lesson:JA-C465-enogu}



\begin{quote}
Before the new one: say the Japanese for to attack, then the Japanese for to snatch.
\end{quote}

\subsection*{You'll want to know: \ja{えのぐ}}
\textbf{\ja{えのぐ}} — \emph{enogu} — \textquotedblleft{}paints\textquotedblright{}.

\begin{quote}
\textbf{\ja{あかい} \ja{えのぐを} \ja{つかいました。}} — \emph{akai enogu o tsukaimashita} — \textquotedblleft{}I used red paint.\textquotedblright{}
\end{quote}

\begin{cousinweb}
\textbf{\ja{えのぐ}} (\emph{enogu}) is a mixed phrase: Sino-Japanese \emph{e}, \textquotedblleft{}picture\textquotedblright{}, and \emph{gu}, \textquotedblleft{}tool\textquotedblright{} (Middle Chinese \emph{gjuH}), both borrowed from Chinese, joined by the native particle \emph{no}, \textquotedblleft{}of\textquotedblright{} (JMdict, KANJIDIC2; Middle Chinese in Baxter's transcription).
\end{cousinweb}

\begin{grammarlens}[title={What you've built}]
One more word for describing what you lived and felt.
\end{grammarlens}

\subsection*{Your turn}
\begin{itemize}
\raggedright
  \item \emph{Say it:} \emph{enogu}
  \item \emph{Say it:} \emph{enogu}, once more
  \item \emph{Say it:} \emph{akai enogu o tsukaimashita} — I used red paint.
  \item \emph{Recall:} say the Japanese for to attack, then the Japanese for to snatch, then say \emph{enogu} again
\end{itemize}

\begin{tcolorbox}[breakable,colback=teal!4,colframe=teal!35!black,title={Before you move on}]
What does \ja{えのぐ} mean? (\textquotedblleft{}Paints\textquotedblright{}.) Say it once more.
\end{tcolorbox}
\section[\texorpdfstring{\ja{ふで} (fude)}{ふで (fude)}]{\ja{ふで} (fude) — a brush; a writing brush}

\label{lesson:JA-C465-fude}



\begin{quote}
Before the new one: say the Japanese for to snatch, then the Japanese for paints.
\end{quote}

\subsection*{You'll want to know: \ja{ふで}}
\textbf{\ja{ふで}} — \emph{fude} — \textquotedblleft{}a brush; a writing brush\textquotedblright{}.

\begin{quote}
\textbf{\ja{ふでで} \ja{なまえを} \ja{かきました。}} — \emph{fude de namae o kakimashita} — \textquotedblleft{}I wrote my name with a brush.\textquotedblright{}
\end{quote}

\begin{cousinweb}
\textbf{\ja{ふで}} (\emph{fude}) is a native Japanese word (wago), inherited rather than borrowed from Chinese; its kanji spelling was attached to it later (JMdict).
\end{cousinweb}

\begin{grammarlens}[title={What you've built}]
One more word for describing what you lived and felt.
\end{grammarlens}

\subsection*{Your turn}
\begin{itemize}
\raggedright
  \item \emph{Say it:} \emph{fude}
  \item \emph{Say it:} \emph{fude}, once more
  \item \emph{Say it:} \emph{fude de namae o kakimashita} — I wrote my name with a brush.
  \item \emph{Recall:} say the Japanese for to snatch, then the Japanese for paints, then say \emph{fude} again
\end{itemize}

\begin{tcolorbox}[breakable,colback=teal!4,colframe=teal!35!black,title={Before you move on}]
What does \ja{ふで} mean? (\textquotedblleft{}A brush; a writing brush\textquotedblright{}.) Say it once more.
\end{tcolorbox}
\section[\texorpdfstring{\ja{えがく} (egaku)}{えがく (egaku)}]{\ja{えがく} (egaku) — to paint a picture; to depict}

\label{lesson:JA-C465-egaku}



\begin{quote}
Before the new one: say the Japanese for paints, then the Japanese for a brush.
\end{quote}

\subsection*{You'll want to know: \ja{えがく}}
\textbf{\ja{えがく}} — \emph{egaku} — \textquotedblleft{}to paint a picture; to depict\textquotedblright{}.

\begin{quote}
\textbf{\ja{かれは} \ja{うみを} \ja{えがきました。}} — \emph{kare wa umi o egakimashita} — \textquotedblleft{}He painted the sea.\textquotedblright{}
\end{quote}

Polite \textbf{\ja{えがきます}}, \emph{egakimasu}; a godan verb. Past (polite): \textbf{\ja{むすめが} \ja{はなを} \ja{えがきました。}}, \emph{musume ga hana o egakimashita}.

\begin{cousinweb}
\textbf{\ja{えがく}} (\emph{egaku}) is a native Japanese word (wago), inherited rather than borrowed from Chinese; its kanji spelling was attached to it later (JMdict).
\end{cousinweb}

\begin{grammarlens}[title={What you've built}]
One more word for describing what you lived and felt.
\end{grammarlens}

\subsection*{Your turn}
\begin{itemize}
\raggedright
  \item \emph{Say it:} \emph{egaku}
  \item \emph{Say it:} \emph{egaku}, once more
  \item \emph{Say it:} \emph{kare wa umi o egakimashita} — He painted the sea.
  \item \emph{Recall:} say the Japanese for paints, then the Japanese for a brush, then say \emph{egaku} again
\end{itemize}

\begin{tcolorbox}[breakable,colback=teal!4,colframe=teal!35!black,title={Before you move on}]
What does \ja{えがく} mean? (\textquotedblleft{}To paint a picture; to depict\textquotedblright{}.) Say it once more.
\end{tcolorbox}
\section[\texorpdfstring{\ja{がか} (gaka)}{がか (gaka)}]{\ja{がか} (gaka) — a painter; an artist}

\label{lesson:JA-C465-gaka}



\begin{quote}
Before the new one: say the Japanese for a brush, then the Japanese for to paint a picture.
\end{quote}

\subsection*{You'll want to know: \ja{がか}}
\textbf{\ja{がか}} — \emph{gaka} — \textquotedblleft{}a painter; an artist\textquotedblright{}.

\begin{quote}
\textbf{\ja{おじさんは} \ja{がかです。}} — \emph{ojisan wa gaka desu} — \textquotedblleft{}My uncle is a painter.\textquotedblright{}
\end{quote}

\begin{cousinweb}
\textbf{\ja{がか}} (\emph{gaka}) is Sino-Japanese (kango): two Chinese characters read with their on'yomi, readings borrowed from Chinese. They are \emph{ga}, \textquotedblleft{}picture\textquotedblright{} (Middle Chinese \emph{hweaH}) and \emph{ka}, \textquotedblleft{}expert\textquotedblright{} (\emph{kae}) (KANJIDIC2; Middle Chinese in Baxter's transcription).
\end{cousinweb}

\begin{grammarlens}[title={What you've built}]
One more word for describing what you lived and felt.
\end{grammarlens}

\subsection*{Your turn}
\begin{itemize}
\raggedright
  \item \emph{Say it:} \emph{gaka}
  \item \emph{Say it:} \emph{gaka}, once more
  \item \emph{Say it:} \emph{ojisan wa gaka desu} — My uncle is a painter.
  \item \emph{Recall:} say the Japanese for a brush, then the Japanese for to paint a picture, then say \emph{gaka} again
\end{itemize}

\begin{tcolorbox}[breakable,colback=teal!4,colframe=teal!35!black,title={Before you move on}]
What does \ja{がか} mean? (\textquotedblleft{}A painter; an artist\textquotedblright{}.) Say it once more.
\end{tcolorbox}
\section[\texorpdfstring{\ja{びじゅつ} (bijutsu)}{びじゅつ (bijutsu)}]{\ja{びじゅつ} (bijutsu) — fine art}

\label{lesson:JA-C465-bijutsu}



\begin{quote}
Before the new one: say the Japanese for to paint a picture, then the Japanese for a painter.
\end{quote}

\subsection*{You'll want to know: \ja{びじゅつ}}
\textbf{\ja{びじゅつ}} — \emph{bijutsu} — \textquotedblleft{}fine art\textquotedblright{}.

\begin{quote}
\textbf{\ja{あねは} \ja{だいがくで} \ja{びじゅつを} \ja{べんきょうしています。}} — \emph{ane wa daigaku de bijutsu o benkyō shite imasu} — \textquotedblleft{}My elder sister studies fine art at university.\textquotedblright{}
\end{quote}

\begin{cousinweb}
\textbf{\ja{びじゅつ}} (\emph{bijutsu}) is Sino-Japanese (kango): two Chinese characters read with their on'yomi, readings borrowed from Chinese. They are \emph{bi}, \textquotedblleft{}beauty\textquotedblright{} (Middle Chinese \emph{mijX}) and \emph{jutsu}, \textquotedblleft{}art\textquotedblright{} (\emph{zywit}) (KANJIDIC2; Middle Chinese in Baxter's transcription).
\end{cousinweb}

\begin{grammarlens}[title={What you've built}]
One more word for describing what you lived and felt.
\end{grammarlens}

\subsection*{Your turn}
\begin{itemize}
\raggedright
  \item \emph{Say it:} \emph{bijutsu}
  \item \emph{Say it:} \emph{bijutsu}, once more
  \item \emph{Say it:} \emph{ane wa daigaku de bijutsu o benkyō shite imasu} — My elder sister studies fine art at university.
  \item \emph{Recall:} say the Japanese for to paint a picture, then the Japanese for a painter, then say \emph{bijutsu} again
\end{itemize}

\begin{tcolorbox}[breakable,colback=teal!4,colframe=teal!35!black,title={Before you move on}]
What does \ja{びじゅつ} mean? (\textquotedblleft{}Fine art\textquotedblright{}.) Say it once more.
\end{tcolorbox}
